% Source for the figure of PathIterator (sec. 1.8): a binary search tree,
% the path pi from the root to the current node (in violet), on the node 3,
% after one step (the node 4, reached by popping 3 and 2) and after another
% (the node 5, reached by pushing 6 and 5). Build: make figures (in ../build).
\documentclass[border=4pt]{standalone}
\usepackage[T1]{fontenc}
\usepackage{lmodern}
\usepackage{amsmath,amssymb}
\usepackage{tikz}
\usepackage{lamalunga}
% tikz-ads.tex: the TikZ styles shared by the figures of the memory of the
% manual, \input by each standalone figure: a cell of a vector (cell), the
% cell that an operation touches (touched), a free cell (free), a node of a
% linked structure (node), the labels of indices (index) and of variables
% (var), a pointer (ptr), and \cells and \indices to draw a row of cells.
\usetikzlibrary{arrows.meta,positioning,calc,fit,shapes.multipart}
\definecolor{adsTouched}{RGB}{214,196,236}
\definecolor{adsFree}{RGB}{236,236,236}
\tikzset{
  ads/.style={>={Stealth[length=2mm]},font=\small,
              every node/.append style={inner sep=0pt}},
  cell/.style={draw,minimum width=7mm,minimum height=7mm,
               font=\small\ttfamily,outer sep=0pt},
  touched/.style={cell,fill=adsTouched},
  free/.style={cell,fill=adsFree,text=black!45},
  node/.style={draw,rounded corners=1pt,minimum width=8mm,
               minimum height=6mm,font=\small\ttfamily},
  index/.style={font=\scriptsize\ttfamily,text=black!60},
  var/.style={font=\small\itshape},
  ptr/.style={->,thick}
}
% \cells{name}{x,y}{list}: a row of cells, the first at (x,y), named
% name-0, name-1, ...; each item is text/style, the style one of cell,
% touched and free (e.g., {a/cell,e/touched,/free})
\newcommand{\cells}[3]{%
  \foreach \v/\s [count=\i from 0] in {#3}
    \node[\s] (#1-\i) at ($(#2)+(\i*7mm,0)$) {\v};}
% \indices{name}{n}: the indices 0 .. n-1 under the cells of name
\newcommand{\indices}[2]{%
  \foreach \i in {0,...,\numexpr#2-1\relax}
    \node[index,below=1.2mm of #1-\i] {\i};}

\begin{document}
\newcommand{\tree}[4]{% x offset, current node, list of path edges, label
  \begin{scope}[xshift=#1]
    \foreach \k/\x/\y in {4/0/0,2/-12/-12,6/12/-12,1/-18/-24,3/-6/-24,5/6/-24,7/18/-24}
      \coordinate (n\k) at (\x mm,\y mm);
    \foreach \a/\b in {4/2,4/6,2/1,2/3,6/5,6/7}
      \draw[edge] (n\a) -- (n\b);
    \foreach \a/\b in {#3}
      \draw[edge-touched] (n\a) -- (n\b);
    \foreach \k in {1,...,7}
      \node[vertex] at (n\k) {\k};
    \node[vertex-touched] at (n#2) {#2};
    \node[var] at (0,-33mm) {#4};
  \end{scope}}
\begin{tikzpicture}[ads]
  \tree{0mm}{3}{4/2,2/3}{$\pi = \langle 4, 2, 3 \rangle$}
  \tree{50mm}{4}{}{$\pi = \langle 4 \rangle$}
  \tree{100mm}{5}{4/6,6/5}{$\pi = \langle 4, 6, 5 \rangle$}
  \draw[ptr] (22mm,-12mm) -- (30mm,-12mm) node[midway,above=1.5mm,var] {\texttt{++}};
  \draw[ptr] (72mm,-12mm) -- (80mm,-12mm) node[midway,above=1.5mm,var] {\texttt{++}};
\end{tikzpicture}
\end{document}
